\hypertarget{cont8:paper-ix-prescribed-parameter-polynomial-lifting}{%
\subsection{Paper IX: prescribed-parameter polynomial
lifting}\label{cont8:paper-ix-prescribed-parameter-polynomial-lifting}}

\textbf{Komorebi P. --- 8 September 2026.}

\hypertarget{cont8:setup-and-exact-lift-set}{%
\subsubsection{Setup and exact lift
set}\label{cont8:setup-and-exact-lift-set}}

Let \(q>0\), \(L\ge1\), and \(m,k,e\ge0\). Write
\(S=\mathbb Z/q\mathbb Z\). Let \[
F\in\mathbb Z[Y_1,\ldots,Y_m,T_1,\ldots,T_k]^e,
\qquad x=(y_0,t_0)\in\mathbb Z^{m+k},
\qquad F(x)\equiv0\pmod{q^L}.
\] The actual lower root supplies the unique integral residual \[
r=q^{-L}F(x)\in\mathbb Z^e,\qquad \bar r=r\bmod q\in S^e.
\] Let the actual reduced Jacobian blocks be \[
J=(D_YF)(x)\bmod q:S^m\to S^e,
\qquad K=(D_TF)(x)\bmod q:S^k\to S^e.
\] For a prescribed parameter digit \(v\in S^k\), choose its canonical
integral representative \(\widehat v\). Define \(\mathcal L_v\) to be
the set of actual residue vectors
\(z\in(\mathbb Z/q^{L+1}\mathbb Z)^{m+k}\) satisfying \[
z\equiv x\pmod{q^L},\qquad F(z)\equiv0\pmod{q^{L+1}},\qquad
z_T\equiv t_0+q^L\widehat v\pmod{q^{L+1}}.
\] The formal predicate uses canonical integer representatives for
evaluation and the lower congruence. The previously checked polynomial
congruence theorem and digit-to-residue equivalence establish the actual
residue-class interpretation. Changing an integral representative of a
parameter digit by a multiple of \(q\) does not change its prescribed
residue modulo \(q^{L+1}\).

\hypertarget{cont8:the-obstruction-theorem}{%
\subsubsection{The obstruction theorem}\label{cont8:the-obstruction-theorem}}

Define the finite solved-block cokernel and its parameter-dependent
class by \[
\mathcal C_x=S^e/\operatorname{im}J,
\qquad o_x(v)=[\bar r+Kv]\in\mathcal C_x.
\] Then \[
\boxed{\mathcal L_v\ne\varnothing\iff o_x(v)=0.}
\] More precisely, the checked digit-to-lift equivalence is \[
\boxed{\mathcal L_v\cong\{u\in S^m:Ju=-\bar r-Kv\}.}
\] It sends a solved digit \(u\) to the actual next residue represented
by \(x+q^L(\widehat u,\widehat v)\). The obstruction depends on the
solved block \(J\), not on the range of the full map \((J\ K)\).

Every nonempty \(\mathcal L_v\) has a checked free and transitive action
of \(\ker J\). The action is defined without choosing a base lift.
Choosing one valid solved digit then gives an explicit bijection
\(\mathcal L_v\cong\ker J\). In particular, for every positive modulus,
\[
o_x(v)=0\quad\Longrightarrow\quad\#\mathcal L_v=\#\ker J.
\] These results permit composite \(q\), the degenerate modulus \(q=1\),
and zero dimensions. Neither \(J\) nor the full reduced Jacobian is
assumed surjective.

\hypertarget{cont8:exact-prime-field-counts}{%
\subsubsection{Exact prime-field
counts}\label{cont8:exact-prime-field-counts}}

When \(q=p\) is prime, put \(\rho=\operatorname{rank}_{\mathbb F_p}J\).
Then \[
\boxed{\#\mathcal L_v=
\begin{cases}
p^{m-\rho},&o_x(v)=0,\\
0,&o_x(v)\ne0,
\end{cases}}
\qquad
\boxed{\#\mathcal C_x=p^{e-\rho}.}
\] The count is of distinct actual residue classes at level \(L+1\). The
solved dimension \(m\), rather than the total number \(m+k\) of
variables, appears in the exponent because the parameter digit is fixed.
These rank formulas are asserted only over prime fields.

The main wrapper assumes only the original lower polynomial congruence.
It constructs one residual \(r\) and supplies the criterion and count
uniformly for every \(v\), together with the exact obstruction
cardinality.

\hypertarget{cont8:proof-and-dependency-chain}{%
\subsubsection{Proof and dependency
chain}\label{cont8:proof-and-dependency-chain}}

The actual derivative calculation gives \[
(DF(x)\bmod q)(u,v)=Ju+Kv.
\] Taylor expansion at depth \(L\ge1\) makes the quadratic remainder
divisible by \(q^{L+1}\). Therefore the digit \((u,v)\) is a polynomial
lifting correction exactly when \[
\bar r+Ju+Kv=0.
\] The quotient class \([\bar r+Kv]\) vanishes exactly when this
equation admits \(u\). This uses additive inverses in the image of
\(J\), so the sign agrees with the actual Taylor equation.

The earlier unrestricted correction-digit equivalence is then restricted
to the predicate that the parameter digit equals \(v\). Taking free
coordinates commutes with forming digit residues, and digit residues are
injective in the digit. Consequently the restricted equivalence enforces
exactly the actual parameter residue condition on \(\mathcal L_v\).

Translation by \(\ker J\) acts freely and transitively on the
solved-digit fiber. Transporting this explicit action through the
residue equivalence yields the torsor on \(\mathcal L_v\). Rank-nullity
over \(\mathbb F_p\) then gives the two displayed cardinalities. No
numerical enumeration is used to establish the general theorem.

\hypertarget{cont8:why-the-solved-block-matters}{%
\subsubsection{Why the solved block
matters}\label{cont8:why-the-solved-block-matters}}

Take \(F(Y,T)=T\), \(q=2\), \(L=1\), and \(x=(0,0)\). The full Jacobian
\((0\ 1)\) is surjective, but \(J=0\). Prescribing \(v=1\) forces
\(T\equiv2\pmod4\), which cannot satisfy \(F=0\pmod4\). Prescribing
\(v=0\) gives exactly two lifts, corresponding to \(Y\equiv0,2\pmod4\).

The Lean controls check this actual polynomial example. Thus
full-Jacobian surjectivity alone does not make every prescribed
coordinate choice admissible. The automatic coordinate selection in the
earlier chart theorem is compatible with this obstruction.
